\section{Base Faithfulness}

Every result in this paper must trace back to the eight base atoms, and to nothing else. This appendix is the audit that confirms it does. It asks one hard question of each finding. Does the finding survive the reality that physical space is the flat cusp of the growing honeycomb, rather than an idealized clean lattice? The honest answer, set out below, is that the program is robust. Most findings never lived on the flat slice at all. The ones that did live there hold on it.

The discipline behind this audit is simple to state and strict to keep. The base is a fixed, small list of atoms. No finding earns a physical name by quietly adding a ninth ingredient. If a phenomenon does not come out of the eight, the appendix reports the honest negative. It does not impose the phenomenon and then call it emergent.

\subsection{The eight atoms, and the rule against a ninth}

The base is exactly eight atoms. They cover the universal furniture of any law, a where, a when, a what, and a how, with two laws under the how, because two things evolve, the state and the space.

\begin{axiom}[The eight base atoms]
The flow $F$ is assembled from exactly eight atoms,
\[
F = \langle\, \text{mesh},\ \text{dock},\ \text{site},\ \text{tone},\ \text{lean},\ \text{knit},\ \text{wake},\ \text{beat}\,\rangle ,
\]
written with their eight distinct first letters $m\, d\, s\, t\, l\, k\, w\, b$.
\end{axiom}

Each atom has one job, and only one. A plain statement of each follows, because the audit only makes sense once the list is fixed.

\begin{enumerate}
\item \textbf{mesh} (where). The $\{3,4,3,4\}$ honeycomb, a fixed discrete graph of docks, open and ever growing.
\item \textbf{dock} (where). The here-now cell, a place of transit with $24$ sites.
\item \textbf{site} (where). One of the $24$ directions out of a dock, the directed slot a tone lives in.
\item \textbf{tone} (what). The ternary value $\{-1, 0, +1\}$ carried on the sites, read as pain, peace, pleasure.
\item \textbf{lean} (what). The value order $-1 < 0 < +1$, which gives charge as the signed sum and gives the create move its direction.
\item \textbf{knit} (how the state changes). The deterministic, reversible, charge-conserving local collision, including the create move, followed by the stream.
\item \textbf{wake} (how the space changes). New docks born at peace at the open, ever-expanding edge.
\item \textbf{beat} (when). Discrete synchronous time, reversible, with no built-in direction.
\end{enumerate}

Two of the most physical-sounding things in the paper are deliberately \emph{not} on this list. The arrow, the direction of time, is emergent. It is the wake acting through the lean and the create move, a low-entropy source biasing an otherwise reversible step. The attraction, binding by radiation, is also emergent. It is the radiation-pressure shadow that the active vacuum casts on the open, growing mesh. Naming either of these a base atom would be the very smuggling this appendix exists to forbid.

\begin{principle}[No ninth ingredient]
No finding is allowed a ninth ingredient. We never add cohesion, planning, a stabilizer, or a gravity term and then call the result emergent. If a phenomenon does not come from the eight, we report the honest negative. The raw substrate is not physics directly. It passes through one to three emergent middle layers before anything earns a physical name.
\end{principle}

This rule has teeth precisely because it can be violated and the violation would be invisible from a distance. A model that wants stable matter can quietly add a stabilizing term. A model that wants binding can quietly add an attraction. The faithfulness claim of this paper is that neither was done. Where stability or binding appears, it was measured to come from the knit and the wake, not inserted by hand. The rest of the appendix tests that claim against the hardest geometric fact the model faces.

\subsection{Three layers, three fates}

Physical space, in this theory, is the flat cusp of the honeycomb, the parabolic point at infinity whose flat cross-section is the everyday $3$D world. It is not a clean infinite lattice. So the audit's organizing question is, which findings depend on the exact shape of that flat slice, and which never touched it at all?

The findings sort into three kinds. The cusp question reaches each kind differently.

\begin{itemize}
\item \textbf{Bulk findings, untouched.} Spin, triality, the gauge group, the weak mixing angle as $3/8$, the mass relations, the holographic correlator, cosmology, the hierarchy, and the toolkit. These live in the $4$D bulk sites or are pure group theory. The shape of physical space does not reach them.
\item \textbf{Boundary findings, untouched.} Holography and quantum nonlocality live on the $3$-sphere boundary. The cusp question does not touch them either.
\item \textbf{Flat-layer findings, re-tested.} Dirac, Lorentz, the inverse-square gravity law, electromagnetism, solitons, the self as a fermion, and matter formation. These need clean $3$D space, so each was re-run on the real cusp rather than on an idealized slice.
\end{itemize}

\begin{result}[Most findings are off the slice]
Only the flat-layer findings depend on the exact shape of physical space. The bulk findings live in the $4$D sites or in group theory, and the boundary findings live on the $3$-sphere. Neither class is touched by the cusp question. The audit therefore narrows to a re-test of the flat layer.
\end{result}

\subsection{The re-test, measured}

The flat-layer observables were re-run on the cubic cusp, the parabolic point whose flat cross-section is physical $3$D space. They were also run on a generic disordered slice as a control. The control is the load-bearing part of the test. A re-test that could only pass would prove nothing. This one could have failed on the control, and the contrast between the two columns is the evidence.

\begin{table}[ht]
\centering
\begin{tabular}{@{}lll@{}}
\toprule
observable & cubic cusp, physical space & generic slice, bad extraction \\
\midrule
spectral dimension & $2.97$, clean $3$D & $2.16$, sub-$3$D \\
isotropy of the light cone & $0.14$ & $0.32$, rough \\
\bottomrule
\end{tabular}
\caption{Flat-layer observables re-run on the real cubic cusp and on a generic disordered slice. The cusp recovers clean $3$D physics. The generic slice degrades.}
\label{tab:cusp-retest}
\end{table}

The reading is direct. The flat-layer physics holds on the cubic cusp and degrades on a generic slice. The spectral dimension lands at $2.97$ on the cusp, a clean $3$, and slumps to $2.16$ on the disordered slice. The light cone is nearly isotropic on the cusp and rough on the slice.

\begin{result}[The flat layer holds on the cusp]
The flat-layer physics holds on the cubic cusp and degrades on a generic disordered slice. The past results therefore survive provided physical space is the cusp cubic, which it is. They would fail on a disordered slice, which is for that reason not physical space. The control could have failed, and the contrast is the evidence.
\end{result}

\subsection{The worst-first problem ledger}

Faithfulness is not the same as having no open problems. It means naming the open ones plainly and not papering over them. What follows is the honest list of the deepest worries on the real cusp, ranked worst first, each with its status.

\begin{table}[ht]
\centering
\begin{tabular}{@{}p{0.42\textwidth}p{0.52\textwidth}@{}}
\toprule
problem & status \\
\midrule
does the $4$D bulk spinor live on the $3$D cusp & solved, it is a projection, the $4$ splits into $2 + 2$, verified \\
cubic anisotropy, the cusp lattice breaks rotations & resolved, tiny, of order the squared ratio of energy to the Planck scale, isotropic in the infrared, falsifiable, smaller than the naive cubic estimate \\
dynamical settling, does matter actually gather into clean cusp regions & open, a plausible mechanism and a defined test, not yet confirmed, the one genuinely open robustness item \\
absolute quantitative values, masses, couplings, the Born-rule normalization & the frontier, shared with every theory of this kind \\
nonlocality and the noncompact geometry & resolved, it is a feature, the cusp is the $3$D space and the nonlocality is bulk-provided \\
\bottomrule
\end{tabular}
\caption{The worst-first robustness ledger on the real cusp.}
\label{tab:ledger}
\end{table}

The ledger has one genuinely open item, dynamical settling. The question is whether matter actually gathers into clean cusp regions on its own, rather than being placed there. There is a plausible mechanism and a defined test. There is not yet a confirmation. The appendix records this as open and does not pretend otherwise.

One subtlety in the ledger deserves a sentence of its own. The cubic anisotropy belongs to the cusp, not to the bulk.

\begin{proposition}[The residual anisotropy is the cusp's]
The bulk sites are isotropic to fourth order, protected by the substrate symmetry. Matter, though, propagates on the cubic cusp, whose exact symmetry is the finite octahedral group, anisotropic at fourth order and not protected by the bulk symmetry. The one residual Lorentz violation therefore lives specifically on the cusp.
\end{proposition}

\subsection{The one modification, and the honest artifact}

After the full re-test, the single change the cusp reality makes to the body of the paper is one effect, and it is small.

\begin{result}[The single cusp modification]
The only change the cusp reality makes is the cusp's small cubic anisotropy, a tiny ultraviolet Lorentz violation of order the squared energy-to-Planck ratio, isotropic in the infrared. Everything else is unaffected or holds.
\end{result}

One artifact must be flagged plainly, because leaving it unflagged would itself be a small infidelity. A finite patch of the cusp reads a gravity exponent near $0.67$. That is a finite-size and screening artifact, not a real deviation from the inverse-square law. The clean inverse-square law is the difference-method result, which removes the finite-size confound, and not the raw patch reading.

\begin{principle}[Report the corrected number]
The gravity exponent of $0.67$ read off a finite patch is a finite-size and screening artifact. The faithful number is the inverse-square law recovered by the difference method. We report the corrected number, never the confounded one.
\end{principle}

\subsection{Why it matters}

The program is robust to the cusp reality. Only the flat-layer physics depends on clean $3$D space, and it holds on the cusp while a generic slice fails. The bulk and boundary findings are untouched by the shape of physical space.

\begin{result}[The faithfulness this appendix certifies]
Every result in the body of this paper rests on the eight base atoms, with no smuggled ninth. The flat-layer findings were re-run on the real cusp and hold. The bulk and boundary findings never depended on the cusp. Where a result is open, the ledger names it open. That is the faithfulness this appendix certifies.
\end{result}

The arrow and the attraction stay emergent, the wake and the knit do the work that lesser models would hand to a ninth atom, and the one honest open item, dynamical settling, is named as open rather than waved away. The audit closes where it began. Eight atoms, and nothing else.
